\documentclass[11pt,a4paper]{article}
\usepackage[a4paper,margin=1.45cm]{geometry}
\usepackage{fontspec}
\usepackage{enumitem}
\usepackage[hidelinks]{hyperref}
\usepackage{tabularx}
\usepackage{array}
\usepackage{xcolor}
\usepackage{titlesec}
\usepackage{fancyhdr}
\pagestyle{fancy}
\fancyhf{}
\renewcommand{\headrulewidth}{0pt}
\cfoot{\thepage}

\setmainfont{TeX Gyre Pagella}
\definecolor{titlecolor}{HTML}{18314A}
\definecolor{textgray}{HTML}{4D5B69}
\titleformat{\section}{\large\bfseries\color{titlecolor}}{}{0pt}{}
\titlespacing*{\section}{0pt}{10pt}{6pt}
\setlist[itemize]{leftmargin=1.5em, itemsep=2pt, topsep=3pt}
\setlength{\parindent}{0pt}
\setlength{\parskip}{3pt}

\newcommand{\entry}[4]{%
  \textbf{#1} \hfill #2\\
  {\color{textgray}#3} \hfill {#4}\\[-2pt]
}

\begin{document}

{\LARGE\bfseries\color{titlecolor} Qun Niu}\\[4pt]
Beijing Institute of Technology \quad Ph.D. Candidate in M.E. (Expected graduation: Dec. 2026)\\
I am currently seeking suitable industry positions or postdoctoral opportunities\\
Email: \href{mailto:nq_eecm@163.com}{nq\_eecm@163.com} \quad Mobile: 15601280520\\
Homepage: \href{https://beta1scat.github.io/academic-cv/}{https://beta1scat.github.io/academic-cv/}

\section*{Profile}
I am a Ph.D. candidate in Mechanical Engineering at Beijing Institute of Technology. My research focuses on robot perception, motion planning, and grasp planning, especially for grasping in cluttered or unstructured scenes. I work on integrating 3D vision, geometric reasoning, kinematic analysis, and task execution into coherent robotic systems. My recent work spans inverse kinematics, hand-eye and tool calibration, point-cloud pose estimation, scene understanding through geometric primitive fitting, and language-guided grasping. In addition to geometry-driven and analytical methods, I also incorporate learning-based methods into robotic tasks and am currently exploring VLA and reinforcement learning. I am experienced with C/C++, Python, PyTorch, and ROS, enjoy challenging problems, and adapt quickly to new methods and tools.

\section*{Education}
\entry{Beijing Institute of Technology}{Sep. 2020 -- Present}{Mechanical Engineering}{Ph.D.}
\entry{Beijing University of Technology}{Sep. 2017 -- Jun. 2020}{Mechanical Engineering}{M.S.}
\entry{Hebei University of Science and Technology}{Sep. 2013 -- Jun. 2017}{Process Equipment and Control Engineering}{B.S.}

\section*{Research Interests}
\begin{itemize}
	\item Robot perception: scene segmentation, point-cloud classification, pose estimation, geometric fitting, and 3D vision system integration.
	\item Kinematics, motion planning, and calibration: inverse kinematics, path planning, collision detection, hand-eye calibration, and tool coordinate calibration.
	\item Vision-guided grasping: hand-eye calibration, tool coordinate calibration, grasp pose generation, and grasp filtering.
	\item Language-guided grasping: combining multimodal object localization, geometric reasoning, and task constraints for grasp planning.
\end{itemize}

\section*{Publications}
\begin{itemize}
	\item \textbf{Niu, Qun}, Zhang, Chuanlin, Zhang, Tianyu, Zhao, Jieliang, Fu, Tie, and Chen, Xuemei. Language-Guided Robot Grasping Based on Basic Geometric Shape Fitting. \textit{Advanced Intelligent Systems}, e202501276, 2025.
	\item \textbf{Niu, Qun}, Zhao, Jieliang, Liang, Lulu, Xing, Jin, Li, Hongkun, and Wang, Ziru. A General Framework for the Analytical Inverse Kinematics Solution of Industrial Robots. \textit{Proceedings of the Institution of Mechanical Engineers, Part C: Journal of Mechanical Engineering Science}, 239(12):4499--4511, 2025.
	\item \textbf{Niu, Qun}, Wang, Ziru, Li, Hongkun, and Zhao, Jieliang. A Parameters Optimization Framework for Pose Estimation Algorithm Based on Point Cloud. \textit{Journal of Physics: Conference Series}, 2746(1):012039, 2024.
	\item \textbf{Niu, Qun}, Zhao, Jieliang, Liang, Lulu, and Xing, Jin. Dynamic Vertical Climbing Mechanism of Chinese Dragon-Li Cats Based on the Linear Inverted Pendulum Model. \textit{Journal of Bionic Engineering}, 20(1):136--145, 2023.
	\item \textbf{Niu, Qun}, et al. A Simultaneous Hand-Eye and Tool Coordinate System Calibration Framework for Robotic Grasping. \textit{Intelligent Service Robotics} (under review).
	\item \textbf{Niu, Qun}, et al. Enhancing Scene Perception: Shape Fitting with Geometric Primitives. \textit{Intelligent Service Robotics} (under review).
\end{itemize}

\section*{Contribution Overview}
\begin{itemize}
	\item Proposed a general analytical inverse kinematics framework that enables industrial robots with the same configuration to share one D-H-parameter-based solver.
	\item Introduced a joint calibration model for hand-eye transformation and tool coordinates to reduce error propagation in robotic grasping.
	\item Built a parameter optimization framework for point-cloud pose estimation to reduce repetitive manual tuning for different objects.
	\item Developed a scene understanding pipeline that combines segmentation, point-cloud classification, and geometric primitive fitting for interpretable perception.
	\item Proposed a language-guided model-free grasping method that integrates multimodal localization, primitive fitting, and task-oriented grasp selection.
	\item Analyzed dynamic vertical climbing in domestic cats with a linear inverted pendulum model and clarified the key mechanism behind dynamic wall traversal without static adhesion.
\end{itemize}

\section*{Patents}
\begin{itemize}
	\item Zhao, Jieliang; Niu, Qun; et al. Adjustable dual-slider ramp experimental platform with tunable height and angle. Authorized invention patent, CN202110431126.4.
	\item Niu, Qun; Zhao, Jieliang; et al. Joint-angle generation method for robots and robot system. Authorized invention patent, CN202310966039.8.
	\item Niu, Qun; Zhao, Jieliang; et al. Vision-guided robotic system and robot calibration method. Authorized invention patent, CN202410192128.6.
	\item Niu, Qun; Zhao, Jieliang; et al. Pose-estimation model configuration method. Pose-estimation method. Under substantive examination, CN202311698845.8.
	\item In addition, he has four collaborative papers on quadruped and wall-climbing robots and six more robot-related authorized or pending patents.
\end{itemize}

\section*{Academic Projects}
\entry{Humanoid Motion Coordination for a Dual-Arm Robot}{Sep. 2017 -- Jun. 2020}{Core member}{Academic project}
Designed a cable-driven seven-DOF manipulator, established its D-H-based kinematic model, proposed an inverse kinematics algorithm for humanoid motion, and developed Cartesian-space and joint-space planning methods for shaft-hole assembly tasks.

\entry{Biomimetic Quadruped Robot with Integrated Joint Actuation}{Sep. 2020 -- Aug. 2022}{Core member}{Academic project}
Worked on integrated joint actuator development inspired by insect leg joints and investigated the mechanism of cat-inspired dynamic vertical climbing. High-speed camera data were used to analyze vertical climbing behavior with the SLIP model, and an elastic actuator with active elasticity and passive restraint was designed.

\entry{Target Detection and Path Planning Algorithms for Robotic Grasping}{Aug. 2022 -- Aug. 2024}{Core member}{Academic project}
Focused on robot algorithms, visual algorithms, and grasp planning for disorderly grasping tasks, including kinematics, path planning, collision detection, hand-eye calibration, pose estimation, grasp pose generation, and grasp filtering.

\section*{Internship Experience}
\entry{Beijing Modi Technology Co., Ltd.}{Apr. 2018 -- Apr. 2019}{R\&D Department}{Internship}
Developed forward and inverse kinematics algorithms, joint-space planning, and Cartesian-space planning for a biomimetic four-axis manipulator based on lever mechanisms. Also worked on interface development and coordinated control for a multi-mode cell-culture bioreactor.

\entry{Institute of Automation, Chinese Academy of Sciences}{Apr. 2019 -- May 2020}{Research Center for Digital Content Technology and Service}{Internship}
Participated in polishing-path planning based on quadric surfaces, robot and force-controlled polishing head programming, on-site debugging, and project delivery. Also worked on multi-axis synchronized PLC control for an automatic screw-locking machine and waterproof housing design for a snake robot.

\entry{Beijing Transfertech Co., Ltd.}{Aug. 2022 -- Aug. 2024}{R\&D Department}{Internship}
Developed and maintained robot kinematics, motion planning, and collision detection algorithms. Contributed to grasp pose generation and filtering, exploratory work on model-free grasping, parameter optimization for pose estimation, hand-eye calibration, and a Blazor-based scene configuration module for robot and environment visualization.

\entry{Beijing Galbot Co., Ltd.}{Aug. 2025 -- Feb. 2026}{Algorithm Engineering}{Internship}
Participates in embodied intelligence algorithm development for indoor mobile manipulation platforms, covering perception, decision-making, and control. Contributes grasping skills for a robotic shop-assistant scenario in the Beijing Galbot Co., Ltd. capsule project using both traditional and learning-based methods on the G1 robot.

\section*{Skills}
\begin{itemize}
	\item Programming and tools: C/C++, Python, PyTorch, ROS, Blazor.
	\item Robot algorithms: kinematics, inverse kinematics, path planning, collision detection, grasp pose generation, and grasp filtering.
	\item Vision and perception: hand-eye calibration, pose estimation, point-cloud processing, scene segmentation, geometric fitting, and 3D vision systems.
	\item Learning-based methods: practical use of learning-based methods in robotic tasks, with current interest in VLA and reinforcement learning.
\end{itemize}

\section*{Personal Strengths}
\begin{itemize}
	\item Enjoys digging into technical problems and sustaining effort on complex tasks.
	\item Motivated by challenging work and willing to take on unfamiliar problems.
	\item Adapts quickly to new research topics, engineering settings, and technical tools.
\end{itemize}

\section*{Honors and Awards}
\begin{itemize}
	\item Outstanding Graduate Student Award, Beijing University of Technology.
	\item First Prize for Graduate Academic Excellence.
	\item Graduate Research Excellence Award.
	\item Second Prize, Hebei Mechanical Design Innovation Competition.
	\item Basketball awards: one municipal-level award and five university-level awards.
\end{itemize}

\end{document}
